\section{Principales aportaciones}
El desarrollo de este proyecto ofrece aportaciones en tres entornos: social y técnico.

\subsection{Social}

Dada la gran importancia cultural del timple, especialmente para la sociedad canaria, la ejecución de este proyecto resulta en un espacio dedicado a la comunidad del timplistas y todas aquellas personas interesadas en el mismo.

La aplicación web pone a disposición de los usuarios un amplio abanico de posibilidades: desde el aprendizaje a través de las páginas informativas, hasta la creación y publicación de canciones, así como el uso de herramientas para facilitar la práctica del instrumento.

Este espacio pone, mas cerca que nunca, el timple en manos de todo aquel que quiera tocarlo.

A nivel social podemos encontrar varias aportaciones.

En primer lugar, la existencia de este espacio promueve la difusión del propio instrumento, ayudando a preservar su legado cultural y a aumentar su popularidad entre nuevas generaciones y diferentes regiones geográficas.

En segundo lugar, la aplicación web ofrece a los usuarios recursos para el aprendizaje y práctica del timple, facilitando la accesibilidad al instrumento y fomentando su estudio.

En tercer lugar, la aplicación web fomenta la creación de contenido musical para el timple, facilitando a los usuarios la publicación y compartición de canciones.

En cuarto y último lugar, aunque ya existen comunidades de timplistas, este espacio digital puede servir como punto de encuentro para todos aquellos interesados en el instrumento, promoviendo la interacción y colaboración entre sus miembros.

\subsection{Técnico}
La principal aportación técnica de este proyecto es la creación de un espacio donde el timple sea el protagonista. Para ello, es necesario el desarrollo de utilidades adaptadas tanto a las singularidades técnicas del instrumento como a su contexto cultural. Dada la escasez de recursos y referencias digitales, la creación de esta aplicación web representa un importante desafío.

Actualmente existen espacios dedicados a la práctica de instrumentos como la guitarra y el ukelele. Estos ofrecen una experiencia adaptada a sus necesidades particulares, contando con grandes comunidades y en constante crecimiento a nivel internacional.

